\section*{Chapter \ref{ch:b1:periodicity} --- Periodicity}

\begin{solution}{exo:b1:periodicity:1}
\ce{Ca} $[\ce{Ar}]\,4\text{s}^2$: \omterm{def:b1:periodicity:table}{period} 4, \omterm{def:b1:periodicity:table}{group} 2, s \omterm{def:b1:periodicity:table}{block}. \ce{Fe}
$[\ce{Ar}]\,3\text{d}^6 4\text{s}^2$: \omterm{def:b1:periodicity:table}{period} 4, \omterm{def:b1:periodicity:table}{group} 8, d \omterm{def:b1:periodicity:table}{block}. \ce{Se}
$[\ce{Ar}]\,3\text{d}^{10} 4\text{s}^2 4\text{p}^4$: \omterm{def:b1:periodicity:table}{period} 4, \omterm{def:b1:periodicity:table}{group} 16, p
\omterm{def:b1:periodicity:table}{block}. \ce{I} $[\ce{Kr}]\,4\text{d}^{10} 5\text{s}^2 5\text{p}^5$: \omterm{def:b1:periodicity:table}{period} 5,
\omterm{def:b1:periodicity:table}{group} 17, p \omterm{def:b1:periodicity:table}{block}.
\end{solution}

\begin{solution}{exo:b1:periodicity:2}
\ce{Na} is larger than \ce{Cl}: same \omterm{def:b1:periodicity:table}{period}, smaller $Z^{*}$. \ce{Na} is
larger than \ce{Na+}, which has lost its 3s shell. \ce{Cl-}
(\qty{181}{pm}) is larger than \ce{F-} (\qty{133}{pm}): one more shell.
\ce{Na+} (\qty{102}{pm}) is larger than \ce{Mg^{2+}} (\qty{72}{pm}): same
ten electrons, nuclear charge 11 against 12.
% ledger: rion:Cl-VI, rion:F-VI, rion:Na+VI, rion:Mg2+VI
\end{solution}

\begin{solution}{exo:b1:periodicity:3}
$\ce{Na} (5.14) < \ce{Al} (5.99) < \ce{Mg} (7.65) < \ce{S} (10.36) <
\ce{P} (10.49) < \ce{Ar} (15.76)$. The general increase along the \omterm{def:b1:periodicity:table}{period}
is broken twice: \ce{Al} below \ce{Mg} (the electron removed is the first
3p electron, higher in energy than 3s) and \ce{S} below \ce{P} (the
electron removed is the first paired 3p electron).
\end{solution}

\begin{solution}{exo:b1:periodicity:4}
The \omterm{def:b1:periodicity:polarisability}{polarisability} measures how easily the electron cloud is deformed by
an electric field. \ce{I-} is more polarisable than \ce{F-}: it is larger
and its outer electrons are farther from the nucleus. \ce{Cl-} (18
electrons held by $Z = 17$) is far more polarisable than \ce{Na+} (10
electrons held by $Z = 11$), as anions are than cations.
\end{solution}

\begin{solution}{exo:b1:periodicity:5}
Ratios: $18.83/5.99 = 3.1$, $28.45/18.83 = 1.5$, $119.99/28.45 = 4.2$. The
jump comes after the third electron: three \omterm{def:b1:quantum-numbers:configuration}{valence electrons}, \omterm{def:b1:periodicity:table}{group} 13,
\omterm{def:b1:periodicity:table}{period} 3: aluminium. It forms \ce{Al^{3+}}, $1\text{s}^2 2\text{s}^2
2\text{p}^6$.
% ledger: ie:Al, ie2:Al, ie3:Al, ie4:Al
\end{solution}

\begin{solution}{exo:b1:periodicity:6}
The added electron of \ce{F-} enters the small 2p shell, where it is
strongly repelled by the seven \omterm{def:b1:quantum-numbers:configuration}{valence electrons} already there; in the
larger 3p shell of chlorine the repulsion is smaller, so more energy is
released. In nitrogen ($2\text{p}^3$, \omorbs{u,u,u}) the extra electron
would have to pair in a half-filled \omterm{def:b1:quantum-numbers:orbital}{subshell}; in magnesium ($3\text{s}^2$)
it would have to start the higher 3p \omterm{def:b1:quantum-numbers:orbital}{subshell}. In neither case is the
anion more stable than the atom and a free electron.
\end{solution}

\begin{solution}{exo:b1:periodicity:7}
$\chi_M(\ce{Na}) = (5.14 + 0.55)/2 = \qty{2.85}{eV}$; $\chi_M(\ce{Cl}) =
(12.97 + 3.61)/2 = \qty{8.29}{eV}$. The difference is large: the electron
pair of a \ce{Na-Cl} bond belongs almost entirely to chlorine, and sodium
chloride is ionic, \ce{Na+ Cl-}.
\end{solution}

\begin{solution}{exo:b1:periodicity:8}
\ce{H-F}: $\Delta\chi = 1.78$, at the border of the ionic region (in
practice a very polar covalent bond), \ce{F^{\delta-}}. \ce{C-H}: 0.35,
nearly non-polar. \ce{Na-Cl}: 2.23, ionic, \ce{Cl^{-}}. \ce{C-Cl}: 0.61,
polar covalent, \ce{C^{\delta+}-Cl^{\delta-}}. \ce{O-H}: 1.24, polar
covalent, \ce{O^{\delta-}-H^{\delta+}}.
\end{solution}

\begin{solution}{exo:b1:periodicity:9}
Down the group the valence electron is in a shell of larger $n$, farther
from the nucleus and screened by more \omterm{def:b1:quantum-numbers:configuration}{core electrons}; it is removed more
easily. Rubidium, one shell further, has a \omterm{def:b1:periodicity:ionisation-energy}{first ionisation energy} below
\qty{4.34}{eV}.
\end{solution}

\begin{solution}{exo:b1:periodicity:10}
Zinc is $[\ce{Ar}]\,3\text{d}^{10} 4\text{s}^2$: its electron comes from a
filled 4s; gallium, $\dots 4\text{s}^2 4\text{p}^1$, loses its single 4p
electron, higher in energy and screened by the full 4s and 3d. Arsenic,
$4\text{p}^3$, loses an electron from a half-filled \omterm{def:b1:quantum-numbers:orbital}{subshell}; selenium,
$4\text{p}^4$, loses the paired electron, pushed up by the repulsion of
its partner: its ionisation energy is slightly lower.
\end{solution}

\begin{solution}{exo:b1:periodicity:11}
$\Delta E = (3.98 - 2.20)^2 = 1.78^2 = \qty{3.17}{eV} = 3.17 \times 96.5 =
\qty{306}{kJ/mol}$. Then $D(\ce{H-F}) = \tfrac12(436 + 158) + 306 = 297 +
306 \approx \qty{603}{kJ/mol}$. The large excess is the extra attraction
between the partial charges $\ce{H^{\delta+}}$ and $\ce{F^{\delta-}}$, due
to the large \omterm{def:b1:periodicity:electronegativity}{electronegativity} difference.
\end{solution}

\begin{solution}{exo:b1:periodicity:12}
Pauling's scale is defined from the energies of bonds A--B: an element
that forms no bond has no Pauling value. Helium, neon and argon form no
stable compounds. Krypton and especially xenon have larger, more
polarisable valence shells and lower ionisation energies
(\qty{14.0}{eV} for krypton against \qty{15.76}{eV} for argon); they are
oxidised by fluorine and oxygen and form compounds such as \ce{XeF2}, so
their bond energies can be measured.
% ledger: ie:Kr, ie:Ar
\end{solution}

\begin{solution}{pb:b1:periodicity:1}
\textbf{1.} Each electron is removed from an ion that is more positive and
smaller than the previous one, which holds its electrons more strongly.
\textbf{2.} $15.04/7.65 = 1.97$; $80.14/15.04 = 5.33$; $109.27/80.14 = 1.36$.
\textbf{3.} The jump comes after two electrons: X has two \omterm{def:b1:quantum-numbers:configuration}{valence electrons},
\omterm{def:b1:periodicity:table}{group} 2; in \omterm{def:b1:periodicity:table}{period} 3 it is magnesium.
\textbf{4.} \ce{Mg^{2+}}, $1\text{s}^2 2\text{s}^2 2\text{p}^6$, the
configuration of neon.
\textbf{5.} A metal: \omterm{def:b1:periodicity:table}{group} 2, on the left of the table, with low first
ionisation energies.
\textbf{6.} The third electron comes from the 2p core, of smaller $n$,
much closer to the nucleus and hardly screened, and is removed from a
doubly charged ion.
% ledger: ie:Mg, ie2:Mg, ie3:Mg, ie4:Mg
\textbf{7.} All have ten electrons: $1\text{s}^2 2\text{s}^2 2\text{p}^6$.
\textbf{8.} The same ten electrons are held by nuclear charges 8, 9, 11,
12, 13: the cloud contracts as $Z$ grows.
\textbf{9.} \ce{O^{2-}}: the largest, with the lowest nuclear charge for
its ten electrons, and an anion.
\textbf{10.} $140/54 = 2.6$.
\textbf{11.} $198.8/2 = \qty{99.4}{pm}$; the anion is $181/99.4 = 1.8$
times larger than the \omterm{def:b1:periodicity:radius}{covalent radius}.
% ledger: re:Cl2, rion:Cl-VI
\textbf{12.} $\ce{Mg^{2+}} < \ce{Na+} < \ce{Cl-}$: the two cations are
isoelectronic (charge 12 against 11), and \ce{Cl-} has a third shell.
\textbf{13.} $\chi_M$ = 10.41 (\ce{F}), 8.29 (\ce{Cl}), 7.59 (\ce{Br}),
7.53 (\ce{O}), 6.26 (\ce{C}), in eV.
\textbf{14.} Mulliken: $\ce{F} > \ce{Cl} > \ce{Br} > \ce{O} > \ce{C}$;
Pauling: $\ce{F} > \ce{O} > \ce{Cl} > \ce{Br} > \ce{C}$. Oxygen moves
from fourth to second.
\textbf{15.} $3.98/10.41 = 0.382$, $3.16/8.29 = 0.381$, $2.96/7.59 =
0.390$: nearly constant, about 0.38; for the halogens the two scales are
proportional.
\textbf{16.} The \omterm{def:b1:periodicity:electron-affinity}{electron affinity} of oxygen is small (\qty{1.44}{eV})
because the added electron must pair in the compact 2p shell; the
Mulliken value of the free atom underrates the pull oxygen exerts in its
bonds, which Pauling's scale, built from bonds, measures directly.
\textbf{17.} Chlorine in \ce{HCl}; fluorine in \ce{ClF}.
\textbf{18.} $\Delta\chi = 0.96$: polar covalent.
\textbf{19.} $D(\ce{H-H}) = 2 \times 218.0 = \qty{436.0}{kJ/mol}$.
\textbf{20.} $D(\ce{Cl-Cl}) = 2 \times 121.3 = \qty{242.6}{kJ/mol}$.
\textbf{21.} $D(\ce{H-Cl}) = 218.0 + 121.3 - (-92.3) = \qty{431.6}{kJ/mol}$.
\textbf{22.} $\Delta E = 431.6 - \tfrac12(436.0 + 242.6) = 431.6 - 339.3 =
\qty{92.3}{kJ/mol}$. In symbols, $D(\ce{H-Cl}) = \Delta_fH(\ce{H}) +
\Delta_fH(\ce{Cl}) - \Delta_fH(\ce{HCl})$ while the mean of the symmetric
bonds is $\Delta_fH(\ce{H}) + \Delta_fH(\ce{Cl})$, so $\Delta E =
-\Delta_fH(\ce{HCl})$.
\textbf{23.} $92.3/96.485 = \qty{0.957}{eV}$.
\textbf{24.} $|\chi(\ce{Cl}) - \chi(\ce{H})| = \sqrt{0.957} =
\textbf{0.98}$, against $3.16 - 2.20 = 0.96$ in the tables: Pauling's
arithmetic, done with modern data, gives the tabulated difference to
within 0.02.
\end{solution}
